% GENERATED FILE. Edit canonical lessons, then run npm run generate:books.
% canonical-source-hash: fnv1a64:15bfd95ba8651345
% canonical-lessons: JA-C39-niwa, JA-C39-ofuro, JA-C39-tatami, JA-C39-tana, JA-C39-hako

\chapter{Garden, Bath, Tatami mat, Shelf, Box}
\label{ch:ja-garden-bath-tatami-mat-shelf-box}
\hlchaptermodality{39}

\begin{chapteropening}
\textbf{By the end of this chapter:} \emph{I can say a garden, a bath, a tatami mat, a shelf and a box.}

Complete the last lesson of chapter 39: i can say a garden, a bath, a tatami mat, a shelf and a box.
\end{chapteropening}

\section[niwa]{\ja{にわ} (niwa) — a garden}

\label{lesson:JA-C39-niwa}



\begin{quote}
Before the new one: say the Japanese for chopsticks, then the Japanese for a kitchen.
\end{quote}

\subsection*{You'll want to know: \ja{にわ}}
\textbf{\ja{にわ}} — \emph{niwa} — \textquotedblleft{}a garden\textquotedblright{}.

\begin{grammarlens}[title={What you've built}]
One more word for this chapter.
\end{grammarlens}

\subsection*{Your turn}
\begin{itemize}
\raggedright
  \item \emph{Say it:} \emph{niwa}
  \item \emph{Say it:} \emph{niwa}, once more
  \item \emph{Say it:} say \emph{niwa}
  \item \emph{Recall:} say the Japanese for chopsticks, then the Japanese for a kitchen, then say \emph{niwa} again
\end{itemize}

\begin{tcolorbox}[breakable,colback=teal!4,colframe=teal!35!black,title={Before you move on}]
What does \ja{にわ} mean? (\textquotedblleft{}A garden\textquotedblright{}.) Say it once more.
\end{tcolorbox}
\section[ofuro]{\ja{おふろ} (ofuro) — a bath}

\label{lesson:JA-C39-ofuro}



\begin{quote}
Before the new one: say the Japanese for a kitchen, then the Japanese for a garden.
\end{quote}

\subsection*{You'll want to know: \ja{おふろ}}
\textbf{\ja{おふろ}} — \emph{ofuro} — \textquotedblleft{}a bath\textquotedblright{}.

\begin{grammarlens}[title={What you've built}]
One more word for this chapter.
\end{grammarlens}

\subsection*{Your turn}
\begin{itemize}
\raggedright
  \item \emph{Say it:} \emph{ofuro}
  \item \emph{Say it:} \emph{ofuro}, once more
  \item \emph{Say it:} say \emph{ofuro}
  \item \emph{Recall:} say the Japanese for a kitchen, then the Japanese for a garden, then say \emph{ofuro} again
\end{itemize}

\begin{tcolorbox}[breakable,colback=teal!4,colframe=teal!35!black,title={Before you move on}]
What does \ja{おふろ} mean? (\textquotedblleft{}A bath\textquotedblright{}.) Say it once more.
\end{tcolorbox}
\section[tatami]{\ja{たたみ} (tatami) — a tatami mat}

\label{lesson:JA-C39-tatami}



\begin{quote}
Before the new one: say the Japanese for a garden, then the Japanese for a bath.
\end{quote}

\subsection*{You'll want to know: \ja{たたみ}}
\textbf{\ja{たたみ}} — \emph{tatami} — \textquotedblleft{}a tatami mat\textquotedblright{}.

\begin{grammarlens}[title={What you've built}]
One more word for this chapter.
\end{grammarlens}

\subsection*{Your turn}
\begin{itemize}
\raggedright
  \item \emph{Say it:} \emph{tatami}
  \item \emph{Say it:} \emph{tatami}, once more
  \item \emph{Say it:} say \emph{tatami}
  \item \emph{Recall:} say the Japanese for a garden, then the Japanese for a bath, then say \emph{tatami} again
\end{itemize}

\begin{tcolorbox}[breakable,colback=teal!4,colframe=teal!35!black,title={Before you move on}]
What does \ja{たたみ} mean? (\textquotedblleft{}A tatami mat\textquotedblright{}.) Say it once more.
\end{tcolorbox}
\section[tana]{\ja{たな} (tana) — a shelf}

\label{lesson:JA-C39-tana}



\begin{quote}
Before the new one: say the Japanese for a bath, then the Japanese for a tatami mat.
\end{quote}

\subsection*{You'll want to know: \ja{たな}}
\textbf{\ja{たな}} — \emph{tana} — \textquotedblleft{}a shelf\textquotedblright{}.

\begin{grammarlens}[title={What you've built}]
One more word for this chapter.
\end{grammarlens}

\subsection*{Your turn}
\begin{itemize}
\raggedright
  \item \emph{Say it:} \emph{tana}
  \item \emph{Say it:} \emph{tana}, once more
  \item \emph{Say it:} say \emph{tana}
  \item \emph{Recall:} say the Japanese for a bath, then the Japanese for a tatami mat, then say \emph{tana} again
\end{itemize}

\begin{tcolorbox}[breakable,colback=teal!4,colframe=teal!35!black,title={Before you move on}]
What does \ja{たな} mean? (\textquotedblleft{}A shelf\textquotedblright{}.) Say it once more.
\end{tcolorbox}
\section[hako]{\ja{はこ} (hako) — a box}

\label{lesson:JA-C39-hako}



\begin{quote}
Before the new one: say the Japanese for a tatami mat, then the Japanese for a shelf.
\end{quote}

\subsection*{You'll want to know: \ja{はこ}}
\textbf{\ja{はこ}} — \emph{hako} — \textquotedblleft{}a box\textquotedblright{}.

\begin{grammarlens}[title={What you've built}]
One more word for this chapter.
\end{grammarlens}

\subsection*{Your turn}
\begin{itemize}
\raggedright
  \item \emph{Say it:} \emph{hako}
  \item \emph{Say it:} \emph{hako}, once more
  \item \emph{Say it:} say \emph{hako}
  \item \emph{Recall:} say the Japanese for a tatami mat, then the Japanese for a shelf, then say \emph{hako} again
\end{itemize}

\begin{tcolorbox}[breakable,colback=teal!4,colframe=teal!35!black,title={Before you move on}]
What does \ja{はこ} mean? (\textquotedblleft{}A box\textquotedblright{}.) Say it once more.
\end{tcolorbox}
